\chapter{Acid-Fast Bacillus (AFB)}

\section*{What Is This Test?}
The acid-fast bacillus (AFB) stain and culture detect mycobacteria and other acid-fast organisms in clinical specimens. Acid-fastness is a property of organisms with high concentrations of mycolic acid (long-chain, waxy 2-alkyl, 3-hydroxy fatty acids) in their cell wall that resists decolorization with acid-alcohol after primary staining. The most clinically significant acid-fast organisms are \textit{Mycobacterium tuberculosis} complex (\textit{M. tuberculosis, M. bovis, M. africanum, M. microti}) and nontuberculous mycobacteria/NTM (\textit{M. avium} complex/MAC, \textit{M. kansasii, M. abscessus, M. fortuitum, M. chelonae, M. marinum, M. ulcerans} and >190 other species). Other partially acid-fast organisms include \textit{Nocardia} species, \textit{Rhodococcus} species, \textit{Tsukamurella}, and \textit{Gordonia}, as well as the cysts of \textit{Cryptosporidium} and \textit{Cystoisospora} in stool. The AFB smear (microscopy) provides rapid presumptive results (same day) and correlates with bacillary load and infectiousness in pulmonary TB. AFB culture (in liquid and solid media) is the gold standard for definitive diagnosis and allows species identification and drug susceptibility testing. AFB NAAT (GeneXpert MTB/RIF Ultra) is described in the Tuberculosis chapter.

\section*{Alternative Names}
Acid-Fast Stain; Ziehl-Neelsen Stain; Kinyoun Stain; Auramine-Rhodamine Stain; AFB Smear; AFB Culture; Mycobacterial Culture; TB Smear; TB Culture; Fluorescent AFB Smear

\section*{Principle}
AFB staining methods exploit the mycolic acid-rich, waxy cell wall of mycobacteria that resists penetration by standard aqueous stains but, once penetrated (by heat or detergent/phenol), retains the dye despite subsequent decolorization with strong acid-alcohol (3\% hydrochloric acid in 95\% ethanol). Ziehl-Neelsen (ZN) stain (hot method): The smear is flooded with concentrated carbol fuchsin (basic fuchsin dissolved in phenol) and heated to steaming for 5 minutes, which melts the mycolic acid and allows the dye to penetrate. The slide is decolorized with 3\% acid-alcohol (HCL in 95\% ethanol) and then counterstained with methylene blue or malachite green. AFB appear as red/pink bacilli against a blue (methylene blue) or green (malachite green) background. Required: at least 300 oil-immersion fields examined before reporting negative. Kinyoun stain (cold method): Uses a higher concentration of phenol and basic fuchsin without heating (detergents increase permeability). Slightly less sensitive than ZN but simpler to perform. Auramine-rhodamine fluorescent stain: The smear is stained with auramine O and rhodamine B (fluorochrome dyes that bind mycolic acid), decolorized with acid-alcohol, and counterstained with potassium permanganate (which quenches background fluorescence). AFB appear as fluorescent yellow-orange bacilli against a dark background. Examined under fluorescence microscopy at lower magnification (200--400x), which allows screening of a larger area (>30 fields) in less time, improving sensitivity. Positive smears should be confirmed by ZN or Kinyoun stain (if they are the initial detection method; if auramine-rhodamine is positive, confirmation by overstaining with ZN is performed). Smear results are reported semi-quantitatively per the WHO/IUATLD scale: Scanty (1--9 AFB/100 high-power fields, reporting exact number), 1+ (10--99 AFB/100 fields), 2+ (1--10 AFB per field in 50 fields), 3+ (>10 AFB per field in 20 fields). Centrifugation (concentration) of specimens before smear preparation increases sensitivity by 10--100 fold (e.g., NALC-NaOH decontamination and centrifugation). Mycobacterial culture: Specimens are decontaminated (NALC-NaOH to kill normal flora while preserving mycobacteria), centrifuged to concentrate, and inoculated into liquid medium (BD BACTEC MGIT 960, detecting oxygen consumption by fluorescence) AND solid medium (Lowenstein-Jensen egg-based medium; Middlebrook 7H10 or 7H11 agar-based medium). Liquid media detect growth faster (10--21 days average), solid media allow colony morphology assessment and detection of mixed cultures. Incubation at 35--37\textdegree{}C for 6--8 weeks. Acid-fastness of colonies is confirmed by ZN or Kinyoun staining. Species identification by MALDI-TOF MS, DNA probe (\textit{M. tuberculosis} complex, MAC, \textit{M. kansasii}, \textit{M. gordonae}), 16S rRNA gene sequencing, or line probe assay.

\section*{History}
Robert Koch identified \textit{Mycobacterium tuberculosis} as the tubercle bacillus in 1882 and developed the original alkaline methylene blue stain. Franz Ziehl (1882) and Friedrich Neelsen (1883) improved the staining method using carbol fuchsin and acid decolorization, creating the Ziehl-Neelsen stain \cite{forbes2018practical}. The cold-staining Kinyoun method was developed in 1915. Fluorescence microscopy with auramine was introduced by Hagemann in 1938 and remains the preferred screening method in high-volume laboratories \cite{forbes2018practical}. The Löwenstein-Jensen egg-based medium was developed in the 1930s. The BACTEC radiometric liquid culture system (460 TB) was introduced in the 1980s and revolutionized TB culture by reducing detection time from 6--8 weeks to 1--3 weeks. The BACTEC MGIT 960 nonradiometric system was introduced in the 1990s and is the current automated standard. WHO endorsed Xpert MTB/RIF for rapid molecular TB diagnosis in 2010 \cite{who2021who}. Despite molecular advances, AFB smear and culture remain essential: smear for rapid assessment of infectiousness and treatment response monitoring; culture for definitive diagnosis, species identification, and drug susceptibility testing.

\section*{Why This Test Is Ordered}
\begin{itemize}
  \item Diagnosis of active pulmonary tuberculosis in patients with persistent cough (>2--3 weeks), hemoptysis, fever, night sweats, weight loss, and chest radiograph abnormalities (apical infiltrates, cavitary lesions, miliary pattern, pleural effusion). Three consecutive early-morning sputum specimens (at least 8 hours apart, with at least one early-morning specimen) should be submitted for AFB smear and culture
  \item Rapid assessment of bacillary load and infectiousness in patients with suspected or confirmed TB. Smear-positivity correlates with transmission risk: smear-positive patients are 5--10 times more infectious than smear-negative, culture-positive patients. Smear positivity is used to determine the duration of airborne isolation
  \item Diagnosis of extrapulmonary TB: AFB smear and culture of appropriate specimens (CSF, pleural fluid, lymph node aspirate/biopsy, urine, tissue biopsy, bone marrow, peritoneal fluid, pericardial fluid)
  \item Monitoring of anti-TB treatment response: Serial sputum AFB smears and cultures are performed at the end of the intensive phase (2 months) and at treatment completion. Conversion from sputum smear/culture positive to negative at 2 months is an important predictor of treatment success
  \item Diagnosis of nontuberculous mycobacterial (NTM) infection in patients with chronic cough, bronchiectasis (particularly \textit{M. avium} complex and \textit{M. abscessus}), nodular lung disease (lady Windermere syndrome), skin and soft tissue infection (\textit{M. marinum} from fish tank exposure; \textit{M. fortuitum, M. chelonae, M. abscessus} after cosmetic surgery, trauma, or injections), and lymphadenitis (particularly cervical lymphadenitis/scrofula from NTM in children)
  \item Diagnosis of \textit{Nocardia} infection in immunocompromised patients with pulmonary, CNS, or cutaneous disease. \textit{Nocardia} is partially acid-fast (modified Kinyoun stain: decolorize with 1\% H$_2$SO$_4$ instead of 3\% HCL alcohol). Nocardia is weakly acid-fast whereas \textit{Mycobacterium} is strongly acid-fast
  \item Screening of patients starting TNF-alpha inhibitor therapy (infliximab, adalimumab, etanercept) for latent TB; though IGRA or TST is the primary screen, AFB smear and culture are indicated if active TB is suspected
\end{itemize}

\section*{Is This a Primary or Secondary Test}
AFB smear is the primary rapid diagnostic test for suspected pulmonary TB. AFB culture is the gold-standard primary test for definitive diagnosis and drug susceptibility testing. NAAT (GeneXpert MTB/RIF Ultra) is an equally primary rapid test that provides both TB detection and rifampin resistance results in 2 hours; WHO recommends NAAT as the initial diagnostic test \cite{who2021who}. Culture remains necessary for phenotypic DST and species identification \cite{lewinsohn2017official}.

\section*{How This Test Is Performed}
Sputum: Three consecutive early-morning specimens, collected in sterile leak-proof containers. The patient should rinse the mouth with water, cough deeply, and expectorate 5--10 mL of purulent material (not saliva). Sputum induction with hypertonic (3\%) saline via nebulizer can be used for patients unable to spontaneously produce sputum. Collection should occur in a negative-pressure airborne infection isolation room or a well-ventilated outdoor area. For children or patients who cannot expectorate, gastric aspirate (early morning, before the patient rises and peristalsis empties the stomach) or bronchoalveolar lavage (BAL) may be obtained. Bronchoscopy is high-risk for TB transmission and must be performed in an airborne infection isolation room with N95 respirator use by all personnel. Extrapulmonary specimens are collected per the specific site (CSF via lumbar puncture, pleural fluid via thoracentesis, lymph node via fine needle aspiration or excisional biopsy). All specimens for mycobacterial culture should be transported at room temperature within 24 hours. Sputum specimens should be refrigerated at 2--8\textdegree{}C if transport delay exceeds 1 hour. Freezing is acceptable for molecular testing but reduces culture viability.

\section*{What Preparation Is Needed}
For sputum collection: rinse mouth with water (not mouthwash). Fasting is not required. Early morning specimens are preferred because secretions pool in the airways overnight. For gastric aspiration: the patient should have fasted for 8--10 hours (typically overnight).

\section*{Contraindications and Precautions}
AFB smear sensitivity is only 50--60\% for culture-positive pulmonary TB. A NEGATIVE AFB SMEAR DOES NOT EXCLUDE TB. Smear sensitivity requires 5,000--10,000 AFB/mL of sputum; below this threshold, smears are falsely negative. Up to 20--30\% of culture-positive TB cases are smear-negative (paucibacillary disease), and smear-negative patients still transmit TB (approximately 17\% of transmission is attributable to smear-negative cases). Therefore, empiric anti-TB therapy should NOT be withheld based on a negative AFB smear alone when clinical and radiographic evidence strongly suggests TB. AFB smear cannot distinguish \textit{M. tuberculosis} from NTM; a positive AFB smear in a patient with bronchiectasis may represent NTM (MAC, \textit{M. abscessus}) rather than TB. NAAT (GeneXpert MTB/RIF Ultra) should be performed on the specimen to distinguish MTB from NTM. False-positive AFB smears are rare but can occur due to: (a) contamination of stain reagents with environmental mycobacteria; (b) cross-contamination of specimens during processing; (c) misinterpretation of stain artifact as AFB. NTM species are ubiquitous in water and soil, and a single positive AFB culture for NTM does not necessarily indicate active disease; clinical and radiographic correlation is essential, and ATS/IDSA diagnostic criteria for NTM lung disease require 2 positive sputum cultures (or 1 positive BAL culture) with compatible symptoms and radiographic findings \cite{griffith2007official, daley2020treatment}. The NALC-NaOH decontamination procedure kills a proportion of mycobacteria, reducing culture sensitivity by 10--30\% compared to direct inoculation (but direct inoculation of non-sterile specimens results in overgrowth of normal flora).

\section*{Risks}
Sputum collection is noninvasive and risk-free. Risks of bronchoscopy, gastric aspiration, fine-needle aspiration, and lumbar puncture relate to the procedure, not the laboratory test. The major public health risk is failure to diagnose smear-positive TB promptly, as the patient remains on open wards (not in airborne isolation), exposing healthcare workers, other patients, and visitors. Airborne infection isolation should be instituted for all patients with suspected pulmonary TB pending smear and NAAT results.

\section*{What Happens After the Test}
AFB smear results are available within 24 hours (STAT for high-risk patients). Positive smears are immediately reported to the ordering provider, infection control (to ensure airborne isolation), and public health. NAAT (GeneXpert) results within 2 hours. Mycobacterial culture: preliminary positive (MGIT flagged) at 1--3 weeks (average 14--21 days for MTB); final negative at 6--8 weeks. All positive AFB cultures with MTB complex identification must be reported to public health. Drug susceptibility testing (phenotypic first-line DST: rifampin, isoniazid, ethambutol, pyrazinamide) is performed on the initial MTB isolate and reported at 4--6 weeks from culture positivity (MGIT DST: 10--14 days; proportion method on solid medium: 3--4 weeks).

\section*{Understanding Results}

\subsection*{Below Normal}
\begin{itemize}
  \item \textbf{AFB smear negative (no acid-fast bacilli seen):} AFB were not visualized after examining the required number of fields (300 fields for ZN/Kinyoun, 30 fields for fluorescent). Does not exclude TB. Sensitivity 50--60\% for culture-positive TB. Paucibacillary disease, extrapulmonary TB, early infection, and pediatric TB are commonly smear-negative. If TB is suspected, NAAT and culture must be performed
  \item \textbf{AFB culture negative at 6--8 weeks:} No mycobacteria isolated. In a patient with a positive AFB smear, a negative culture could indicate: (a) nonviable organisms (patient already on effective anti-TB therapy); (b) NTM that failed to grow on the particular media; (c) technical failure. In a smear-negative patient, a negative culture at 6--8 weeks makes TB unlikely
\end{itemize}

\subsection*{Above Normal}
\begin{itemize}
  \item \textbf{AFB smear positive (reported as scanty, 1+, 2+, 3+, 4+):} Acid-fast bacilli visualized. A positive AFB smear on a sputum specimen from a patient with compatible clinical and radiographic findings is highly suggestive of pulmonary TB (PPV >95\% in high-prevalence settings). The semi-quantitative grade correlates with bacillary load and infectiousness: 1+ smears are associated with a lower transmission risk than 3+ or 4+ smears. The patient must remain in airborne isolation until 3 consecutive sputum AFB smears are negative (typically after 2 weeks of effective therapy)
  \item \textbf{AFB culture positive for \textit{M. tuberculosis} complex:} Definitive diagnosis of TB. Species identification (MALDI-TOF, DNA probe) confirms MTB complex vs. NTM. Drug susceptibility testing (DST) should be performed on all initial MTB isolates. First-line DST: rifampin, isoniazid, ethambutol, pyrazinamide. Rifampin resistance +/- isoniazid resistance = MDR-TB. Second-line DST for MDR-TB: fluoroquinolones (moxifloxacin, levofloxacin), injectables (amikacin, kanamycin, capreomycin), and newer/repurposed agents (bedaquiline, linezolid, clofazimine, delamanid)
  \item \textbf{AFB culture positive for nontuberculous mycobacteria (NTM):} Identified to species level. \textit{M. avium} complex (MAC = \textit{M. avium, M. intracellulare, M. chimaera}) is the most common NTM causing pulmonary disease. \textit{M. abscessus} complex (\textit{M. abscessus} subsp. \textit{abscessus}, \textit{massiliense}, and \textit{bolletii}) is the most drug-resistant and difficult to treat. Diagnostic criteria per ATS/IDSA must be met for NTM pulmonary disease: (a) compatible symptoms and radiographic findings (nodular/bronchiectatic or cavitary); plus (b) 2 positive sputum cultures OR 1 positive BAL culture OR a positive lung biopsy with mycobacterial histopathologic features and positive culture. Susceptibility testing for NTM: MAC: clarithromycin (macrolide susceptibility is critical and predicts treatment response); \textit{M. kansasii}: rifampin; \textit{M. abscessus}: clarithromycin (inducible resistance due to \textit{erm}(41) gene), amikacin, cefoxitin, imipenem, tigecycline, linezolid, clofazimine
\end{itemize}

\section*{Normal Results}
\textbf{No acid-fast bacilli seen} on AFB smear. \textbf{No mycobacteria isolated} at 6--8 weeks of culture. For NTM sputum: \textbf{Single positive NTM culture without compatible clinical/radiographic findings:} considered colonization; surveillance, not treatment, is warranted.

\section*{Follow-up Tests}
\begin{itemize}
  \item \textbf{Mycobacterial NAAT (Xpert MTB/RIF Ultra):} Should be performed on the initial smear-positive or smear-negative specimen to rapidly differentiate MTB complex from NTM, detect \textit{M. tuberculosis}, and assess rifampin resistance
  \item \textbf{Mycobacterial culture with drug susceptibility testing:} Phenotypic DST for all initial MTB culture-positive cases. First-line panel: rifampin, isoniazid, ethambutol, pyrazinamide. MDR-TB confirmation: second-line DST for fluoroquinolones, injectables, bedaquiline, linezolid, clofazimine, delamanid
  \item \textbf{Molecular drug resistance testing (line probe assay, GenoType MTBDRplus/MTBDRsl):} Molecular detection of mutations associated with resistance to rifampin, isoniazid (\textit{katG, inhA}), fluoroquinolones (\textit{gyrA, gyrB}), and second-line injectables (\textit{rrs, eis}). Results in 1--2 days from culture or directly from smear-positive sputum
  \item \textbf{Chest radiograph (PA and lateral) or chest CT:} For all patients with positive AFB smear or culture to assess for pulmonary TB or NTM lung disease
  \item \textbf{HIV testing:} Recommended for all patients with TB or NTM infection
  \item \textbf{Baseline and serial liver function tests (ALT, AST, bilirubin):} For patients on anti-TB therapy (isoniazid, rifampin, pyrazinamide), monitor monthly for hepatotoxicity
\end{itemize}

\section*{References}
\bibliographystyle{unsrtnat}
\bibliography{references}

\clearpage
